\documentclass[10pt]{beamer}

% Tema Limpo e Profissional
\usetheme{Madrid}
\usecolortheme{whale}

% Pacotes Essenciais
\usepackage[utf8]{inputenc}
\usepackage[portuguese]{babel}
\usepackage{graphicx}
\usepackage{booktabs}
\usepackage{amsmath}
\usepackage{amssymb}
\usepackage{listings}
\usepackage{xcolor}

% Configuracoes de Codigo (Python)
\lstset{
    language=Python,
    basicstyle=\ttfamily\scriptsize,
    keywordstyle=\color{blue}\bfseries,
    stringstyle=\color{orange},
    commentstyle=\color{gray!80!black}\itshape,
    breaklines=true,
    showstringspaces=false,
    backgroundcolor=\color{gray!5},
    frame=single,
    rulecolor=\color{gray!30},
    numbers=left,
    numberstyle=\tiny\color{gray},
    stepnumber=1,
    numbersep=5pt
}

% Metadados da Apresentacao
\title[Revisão de Programação]{Maratona de Revisão: Lógica e Programação}
\subtitle{25 Questões Práticas: De Operadores Básicos a Funções}
\author[Prof. Reginaldo]{Prof. Reginaldo Fernandes\texorpdfstring{\\ \small{Lógica de Programação}}{}}
\institute[IFCE]{Instituto Federal do Ceará -- Campus Tauá \\ Curso Técnico em Redes de Computadores}
\date{\today}

\begin{document}

% Slide de Titulo
\begin{frame}
    \titlepage
\end{frame}

% Visao Geral e Metodologia
\begin{frame}{Como Utilizar esta Revisão}
    \begin{block}{Objetivo Pedagógico}
        Consolidar a capacidade analítica e algorítmica por meio de \textbf{25 desafios práticos progressivos}.
    \end{block}

    \vspace{0.2cm}
    \begin{columns}
        \begin{column}{0.48\textwidth}
            \textbf{Distribuição dos Tópicos:}
            \begin{itemize}
                \item \textbf{Parte 1 (Q01-Q05):} Operadores Matemáticos, Relacionais e Lógicos.
                \item \textbf{Parte 2 (Q06-Q10):} Estruturas de Decisão (\texttt{if-elif-else}).
                \item \textbf{Parte 3 (Q11-Q15):} Laços de Repetição e Vetores.
                \item \textbf{Parte 4 (Q16-Q25):} Funções (10 questões cobrindo todos os tipos).
            \end{itemize}
        \end{column}
        \begin{column}{0.48\textwidth}
            \textbf{Estrutura de Cada Slide:}
            \begin{itemize}
                \item \textbf{Quesito:} Enunciado claro do problema computacional e especificação de entrada/saída.
                \item \textbf{Ajuda de Resolução:} Guia de raciocínio, dicas de sintaxe e operadores recomendados, \textbf{sem entregar o código final}.
            \end{itemize}
        \end{column}
    \end{columns}
\end{frame}

% Sumario
\begin{frame}{Sumário da Apresentação}
    \tableofcontents
\end{frame}

% Blocos Modulares de Slides
\section{Parte 1: Programação Estruturada e Operadores}

% Slide de Abertura da Secao 1
\begin{frame}
    \begin{center}
        \LARGE \textbf{Parte 1: Programação Estruturada}\\
        \vspace{0.3cm}
        \normalsize \textit{Operadores Aritméticos, Relacionais e Lógicos}\\
        \vspace{0.4cm}
        \small Questões 01 a 05
    \end{center}
    \vspace{0.3cm}
    \begin{alertblock}{Conceito Central}
        A programação estruturada baseia-se no fluxo sequencial e na manipulação de dados na memória através de expressões matemáticas e booleanas.
    \end{alertblock}
\end{frame}

% Questao 1
\begin{frame}{Questão 01: Operadores Aritméticos Básicos}
    \begin{block}{Quesito: Custo de Deslocamento e Consumo}
        Desenvolva um programa que receba três valores:
        \begin{enumerate}
            \item Distância total da viagem em quilômetros (\texttt{km});
            \item Consumo médio do veículo em quilômetros por litro (\texttt{km/l});
            \item Preço do litro de combustível em reais (\texttt{R\$}).
        \end{enumerate}
        O programa deve calcular e exibir:
        \begin{itemize}
            \item O total de litros consumidos ($\text{Distância} / \text{Consumo}$);
            \item O custo total em reais formatado com 2 casas decimais.
        \end{itemize}
    \end{block}

    \begin{exampleblock}{Ajuda de Resolução (Como Pensar)}
        \begin{itemize}
            \item Converta os dados lidos do teclado para \texttt{float}.
            \item Calcule os litros com o operador de divisão (\texttt{/}).
            \item Obtenha o custo multiplicando litros pelo preço unitário (\texttt{*}).
            \item Exiba o resultado com \texttt{f"\{custo:.2f\}"}.
        \end{itemize}
    \end{exampleblock}
\end{frame}

% Questao 2
\begin{frame}{Questão 02: Divisão Inteira e Operador Módulo}
    \begin{block}{Quesito: Conversão de Segundos para H:M:S}
        Um monitor de infraestrutura registra o tempo de atividade de um servidor em segundos inteiros. \\
        Escreva um programa que leia esse valor inteiro e o decomponha no formato: \textbf{Horas}, \textbf{Minutos} e \textbf{Segundos restantes}.
    \end{block}

    \begin{exampleblock}{Ajuda de Resolução (Como Pensar)}
        \begin{itemize}
            \item Lembre-se: $1\text{ h} = 3600\text{ s}$ e $1\text{ min} = 60\text{ s}$.
            \item Calcule as horas completas com divisão inteira: $\text{horas} = \text{total} // 3600$.
            \item Extraia o que sobrou com o operador resto (\texttt{\%}): $\text{resto} = \text{total} \% 3600$.
            \item Obtenha os minutos com $\text{minutos} = \text{resto} // 60$ e os segundos finais com $\text{segundos} = \text{resto} \% 60$.
        \end{itemize}
    \end{exampleblock}
\end{frame}

% Questao 3
\begin{frame}{Questão 03: Operadores Relacionais (Comparação)}
    \begin{block}{Quesito: Monitor de Capacidade de Link de Rede}
        Um link opera a $1000\text{ Mbps}$. Leia o tráfego atual (\texttt{Mbps}) e gere 3 variáveis booleanas (\texttt{True/False}), sem usar \texttt{if}:
        \begin{itemize}
            \item \textbf{Uso Seguro:} Tráfego estritamente menor que $80\%$ da capacidade;
            \item \textbf{Atenção Crítica:} Tráfego atingiu ou superou $90\%$ da capacidade;
            \item \textbf{Sobrecarga:} Tráfego estritamente superior a $100\%$ da capacidade.
        \end{itemize}
        Exiba o resultado booleano de cada uma das três avaliações.
    \end{block}

    \begin{exampleblock}{Ajuda de Resolução (Como Pensar)}
        \begin{itemize}
            \item Operadores de comparação (\texttt{<}, \texttt{>=}, \texttt{>}) retornam diretamente \texttt{True} ou \texttt{False}.
            \item Calcule os limites multiplicando a capacidade por $0.80$, $0.90$ e $1.0$.
            \item Guarde cada comparação direto em uma variável booleana e a imprima.
        \end{itemize}
    \end{exampleblock}
\end{frame}

% Questao 4
\begin{frame}{Questão 04: Operadores Lógicos Comportamentais}
    \begin{block}{Quesito: Controle de Acesso ao Datacenter}
        A catraca eletrônica deve autorizar a abertura da porta quando:
        \begin{itemize}
            \item O usuário possui credencial (\texttt{True/False}) \textbf{E} o horário atual está dentro da janela permitida ($8\text{h}$ às $18\text{h}$);
            \item \textbf{OU} o alarme de emergência está ativado (\texttt{True/False}).
        \end{itemize}
        Construa a expressão lógica em Python que determine se a porta deve destravar.
    \end{block}

    \begin{exampleblock}{Ajuda de Resolução (Como Pensar)}
        \begin{itemize}
            \item Identifique: \texttt{tem\_cracha} (bool), \texttt{hora} (int) e \texttt{emergencia} (bool).
            \item A janela de horário é: \texttt{(hora >= 8 and hora <= 18)}.
            \item Conecte as condições com os operadores \texttt{and} e \texttt{or}.
            \item Use parênteses para garantir que o \texttt{or} avalie a emergência como exceção geral.
        \end{itemize}
    \end{exampleblock}
\end{frame}

% Questao 5
\begin{frame}{Questão 05: Expressões Aritmético-Lógicas}
    \begin{block}{Quesito: Verificador de Regra de Ano Bissexto}
        Um ano é bissexto se atender às regras do calendário gregoriano:
        \begin{itemize}
            \item É divisível por $4$ \textbf{e} não é divisível por $100$; \textbf{OU}
            \item É divisível por $400$.
        \end{itemize}
        Solicite um ano inteiro e avalie se ele é bissexto usando \textbf{uma única expressão booleana} (sem \texttt{if/else}). Imprima apenas o resultado booleano.
    \end{block}

    \begin{exampleblock}{Ajuda de Resolução (Como Pensar)}
        \begin{itemize}
            \item Divisibilidade por $N$ é avaliada com operador de resto: \texttt{ano \% N == 0}.
            \item Não ser divisível por $100$: \texttt{ano \% 100 != 0}.
            \item Combine as condições em uma só linha usando operadores \texttt{and} e \texttt{or} com parênteses adequados.
        \end{itemize}
    \end{exampleblock}
\end{frame}

\section{Parte 2: Estruturas de Decisão}

% Slide de Abertura da Secao 2
\begin{frame}
    \begin{center}
        \LARGE \textbf{Parte 2: Estruturas de Decisão}\\
        \vspace{0.3cm}
        \normalsize \textit{Desvios Condicionais Simples, Compostos e Encadeados}\\
        \vspace{0.4cm}
        \small Questões 06 a 10
    \end{center}
    \vspace{0.3cm}
    \begin{alertblock}{Conceito Central}
        As estruturas de decisão permitem que o fluxo de execução tome caminhos distintos a depender da avaliação de condições lógicas em tempo de execução.
    \end{alertblock}
\end{frame}

% Questao 6
\begin{frame}{Questão 06: Decisão Simples e Composta}
    \begin{block}{Quesito: Radar de Trânsito Inteligente}
        Leia a velocidade registrada de um veículo em via com limite de $60\text{ km/h}$.
        \begin{itemize}
            \item Até $60\text{ km/h}$: exiba \texttt{"Velocidade permitida"}.
            \item Acima de $60\text{ km/h}$: informe \texttt{"Multado!"}, calcule e exiba a multa a $\text{R\$ } 7{,}00$ por $\text{km/h}$ excedente.
        \end{itemize}
    \end{block}

    \begin{exampleblock}{Ajuda de Resolução (Como Pensar)}
        \begin{itemize}
            \item Estruture com \texttt{if-else} testando: \texttt{velocidade > 60}.
            \item No bloco do \texttt{if}, calcule o excesso: $\text{excesso} = \text{velocidade} - 60$.
            \item Multiplique o excesso por $7.00$ para obter o valor total da penalidade.
        \end{itemize}
    \end{exampleblock}
\end{frame}

% Questao 7
\begin{frame}{Questão 07: Decisão Encadeada (\texttt{if-elif-else})}
    \begin{block}{Quesito: Triagem de Glicemia em Jejum}
        Receba a glicose no sangue de um paciente ($\text{mg/dL}$) e classifique:
        \begin{itemize}
            \item Menor que $70\text{ mg/dL}$: \textbf{"Hipoglicemia"}
            \item De $70$ a $99\text{ mg/dL}$: \textbf{"Normal"}
            \item De $100$ a $125\text{ mg/dL}$: \textbf{"Pré-diabetes"}
            \item $126\text{ mg/dL}$ ou mais: \textbf{"Diabetes Estabelecido"}
        \end{itemize}
    \end{block}

    \begin{exampleblock}{Ajuda de Resolução (Como Pensar)}
        \begin{itemize}
            \item Use uma cadeia ordenada de \texttt{if-elif-else} do menor para o maior.
            \item Se o primeiro teste \texttt{glicose < 70} falhar, o próximo ramo já sabe que o valor é $\ge 70$, bastando testar \texttt{elif glicose <= 99:}.
        \end{itemize}
    \end{exampleblock}
\end{frame}

% Questao 8
\begin{frame}{Questão 08: Condições Múltiplas e Geometria}
    \begin{block}{Quesito: Validação e Classificação de Triângulos}
        Receba três lados $A$, $B$ e $C$. O programa deve:
        \begin{enumerate}
            \item Verificar se formam triângulo (cada lado menor que a soma dos outros dois);
            \item Se formarem, classificar: \textbf{Equilátero} (3 iguais), \textbf{Isósceles} (2 iguais) ou \textbf{Escaleno} (todos diferentes);
            \item Caso contrário, avisar que não formam um triângulo.
        \end{enumerate}
    \end{block}

    \begin{exampleblock}{Ajuda de Resolução (Como Pensar)}
        \begin{itemize}
            \item Valide a existência com \texttt{(A < B + C) and (B < A + C) and (C < A + B)}.
            \item Aninhe a classificação dentro do bloco verdadeiro:
            \begin{itemize}
                \item Teste 3 lados iguais: \texttt{A == B == C} (Equilátero).
                \item Teste 2 lados iguais: \texttt{(A == B) or (B == C) or (A == C)} (Isósceles).
                \item O restante cai no \texttt{else} (Escaleno).
            \end{itemize}
        \end{itemize}
    \end{exampleblock}
\end{frame}

% Questao 9
\begin{frame}{Questão 09: Tabela Progressiva de Descontos}
    \begin{block}{Quesito: Política Comercial de Descontos}
        Um e-commerce aplica descontos conforme o valor bruto da compra:
        \begin{itemize}
            \item Até $\text{R\$ } 100{,}00$: $0\%$ | De $\text{R\$ } 100{,}01$ a $\text{R\$ } 300{,}00$: $10\%$
            \item De $\text{R\$ } 300{,}01$ a $\text{R\$ } 500{,}00$: $15\%$ | Acima de $\text{R\$ } 500{,}00$: $20\%$
        \end{itemize}
        Exiba: taxa percentual, valor descontado em $\text{R\$}$ e valor final a pagar.
    \end{block}

    \begin{exampleblock}{Ajuda de Resolução (Como Pensar)}
        \begin{itemize}
            \item Defina uma variável \texttt{taxa} e use \texttt{if-elif-else} apenas para ajustá-la ($0.0$, $0.10$, $0.15$ ou $0.20$).
            \item Fora do bloco condicional, faça os cálculos apenas uma vez: $\text{desconto} = \text{valor} \times \text{taxa}$ e $\text{total} = \text{valor} - \text{desconto}$.
        \end{itemize}
    \end{exampleblock}
\end{frame}

% Questao 10
\begin{frame}{Questão 10: Menu de Operações e Prevenção de Falhas}
    \begin{block}{Quesito: Mini Calculadora com Salvaguarda}
        Apresente um menu (\texttt{1-Soma}, \texttt{2-Subtração}, \texttt{3-Multiplicação}, \texttt{4-Divisão}), receba dois números e execute a opção escolhida.
        \begin{itemize}
            \item Na divisão, evite falha no programa impedindo a divisão por zero.
            \item Para opções inexistentes, exiba \texttt{"Opção Inválida"}.
        \end{itemize}
    \end{block}

    \begin{exampleblock}{Ajuda de Resolução (Como Pensar)}
        \begin{itemize}
            \item Estruture as escolhas usando \texttt{if-elif-else}.
            \item Na opção de divisão, verifique se o segundo número é diferente de zero (\texttt{num2 != 0}) antes de calcular.
            \item Use o \texttt{else} final para tratar qualquer opção digitada fora do intervalo $1$ a $4$.
        \end{itemize}
    \end{exampleblock}
\end{frame}

\section{Parte 3: Estruturas de Repetição e Vetores}

% Slide de Abertura da Secao 3
\begin{frame}
    \begin{center}
        \LARGE \textbf{Parte 3: Repetição e Vetores}\\
        \vspace{0.3cm}
        \normalsize \textit{Laços de Repetição (\texttt{for}, \texttt{while}) e Manipulação de Listas}\\
        \vspace{0.4cm}
        \small Questões 11 a 15
    \end{center}
    \vspace{0.3cm}
    \begin{alertblock}{Conceito Central}
        Laços automatizam o processamento sequencial de dados, viabilizando cálculos estatísticos, buscas lineares e filtros em coleções.
    \end{alertblock}
\end{frame}

% Questao 11
\begin{frame}{Questão 11: Laço Contado e Vetor Acumulador}
    \begin{block}{Quesito: Análise Térmica Semanal}
        Leia a temperatura máxima dos $7$ dias da semana e guarde em uma lista. Ao final, determine e mostre:
        \begin{enumerate}
            \item A média aritmética simples das temperaturas;
            \item A maior temperatura registrada;
            \item Quantos dias tiveram temperatura estritamente superior à média.
        \end{enumerate}
    \end{block}

    \begin{exampleblock}{Ajuda de Resolução (Como Pensar)}
        \begin{itemize}
            \item Use um laço \texttt{for} para coletar as 7 entradas com \texttt{.append()}.
            \item Obtenha a média com $\text{soma} / 7$ (usando \texttt{sum()} ou acumulador).
            \item Em um segundo laço, conte quantos valores na lista são maiores que a média calculada usando uma condição \texttt{if temp > media:}.
        \end{itemize}
    \end{exampleblock}
\end{frame}

% Questao 12
\begin{frame}{Questão 12: Repetição com Sentinela (\texttt{while})}
    \begin{block}{Quesito: Totalizador de Terminal de Saque}
        Simule o fechamento de um caixa de autoatendimento bancário:
        \begin{itemize}
            \item Leia saques contínuos até que o usuário digite $0$ (parada).
            \item Rejeite valores negativos com aviso e desconsidere-os da contagem.
            \item Exiba o total de saques válidos efetuados e a soma total sacada.
        \end{itemize}
    \end{block}

    \begin{exampleblock}{Ajuda de Resolução (Como Pensar)}
        \begin{itemize}
            \item Use um laço \texttt{while True:} com leitura do valor a cada volta.
            \item Se o valor for $0$, encerre a repetição com \texttt{break}.
            \item Se for negativo, alerte o usuário e volte ao topo com \texttt{continue}.
            \item Caso contrário, some o saque ao total e incremente o contador.
        \end{itemize}
    \end{exampleblock}
\end{frame}

% Questao 13
\begin{frame}{Questão 13: Varredura de Vetor e Filtragem}
    \begin{block}{Quesito: Separação de Números Pares e Ímpares}
        Leia $10$ números inteiros para uma lista inicial. Em seguida:
        \begin{itemize}
            \item Crie uma lista com os números pares e outra com os ímpares;
            \item Exiba o tamanho e os elementos contidos nas três listas.
        \end{itemize}
    \end{block}

    \begin{exampleblock}{Ajuda de Resolução (Como Pensar)}
        \begin{itemize}
            \item Crie duas listas vazias: \texttt{pares = []} e \texttt{impares = []}.
            \item Itere na lista com \texttt{for n in lista:}.
            \item Verifique a paridade com \texttt{n \% 2 == 0}: adicione em \texttt{pares} ou em \texttt{impares} com \texttt{.append()}.
            \item Exiba o número de itens com \texttt{len()}.
        \end{itemize}
    \end{exampleblock}
\end{frame}

% Questao 14
\begin{frame}{Questão 14: Busca Linear em Vetores com Flag}
    \begin{block}{Quesito: Verificador de Matrícula em Lista de Acesso}
        Dada uma lista com $8$ matrículas pré-cadastradas, leia uma matrícula informada pelo usuário e faça a busca:
        \begin{itemize}
            \item Se presente, mostre: \texttt{"Autorizado na posição X"}.
            \item Se ausente, mostre: \texttt{"Acesso Negado: Matrícula não cadastrada"}.
        \end{itemize}
    \end{block}

    \begin{exampleblock}{Ajuda de Resolução (Como Pensar)}
        \begin{itemize}
            \item Inicie uma flag booleana de controle: \texttt{achou = False}.
            \item Percorra os índices com \texttt{for i in range(len(lista)):}.
            \item Se o item do índice for igual à busca, marque \texttt{achou = True}, exiba a posição \texttt{i} e interrompa com \texttt{break}.
            \item Ao final do laço, teste: se \texttt{not achou:}, avise que não encontrou.
        \end{itemize}
    \end{exampleblock}
\end{frame}

% Questao 15
\begin{frame}{Questão 15: Inversão Algorítmica de Vetor}
    \begin{block}{Quesito: Espelhamento de Elementos sem Métodos Prontos}
        Preencha uma lista com $6$ números fornecidos pelo usuário. \\
        Crie uma nova lista com os mesmos elementos em \textbf{ordem invertida}, sem usar \texttt{.reverse()} ou fatiamento \texttt{[::-1]}.
    \end{block}

    \begin{exampleblock}{Ajuda de Resolução (Como Pensar)}
        \begin{itemize}
            \item O último elemento está em $\text{tamanho} - 1$ e o primeiro em $0$.
            \item Crie uma nova lista vazia para armazenar o resultado.
            \item Percorra a lista de trás para frente usando: \\
            \texttt{range(len(lista) - 1, -1, -1)}.
            \item Insira cada elemento lido na nova lista com \texttt{.append()}.
        \end{itemize}
    \end{exampleblock}
\end{frame}

\section{Parte 4: Funções e Modularização}

% Slide de Abertura da Secao 4
\begin{frame}
    \begin{center}
        \LARGE \textbf{Parte 4: Modularização com Funções}\\
        \vspace{0.3cm}
        \normalsize \textit{Parâmetros, Retornos e Escopo de Variáveis}\\
        \vspace{0.4cm}
        \small Questões 16 a 25 (10 Desafios Didáticos)
    \end{center}
    \vspace{0.3cm}
    \begin{alertblock}{Conceito Central}
        Funções dividem problemas complexos em sub-rotinas reutilizáveis. Variam em dois eixos: se recebem entradas (\textbf{parâmetros}) e se devolvem respostas (\textbf{\texttt{return}}).
    \end{alertblock}
\end{frame}

% Questao 16
\begin{frame}{Questão 16: Função Sem Parâmetro e Sem Retorno}
    \begin{block}{Quesito: Cabeçalho Padrão Institucional}
        Crie a função \texttt{exibir\_cabecalho()} que imprima na tela um banner formatado com:
        \begin{itemize}
            \item Instituição: \textbf{IFCE Campus Tauá};
            \item Curso: \textbf{Técnico em Redes de Computadores};
            \item Molduras decorativas com o caractere \texttt{"="}.
        \end{itemize}
        Invoque essa função no início do programa principal.
    \end{block}

    \begin{exampleblock}{Ajuda de Resolução (Como Pensar)}
        \begin{itemize}
            \item Declare com \texttt{def exibir\_cabecalho():} (parênteses vazios).
            \item Não utilize a instrução \texttt{return}; utilize comandos \texttt{print()}.
            \item Chame a função no código principal apenas escrevendo: \texttt{exibir\_cabecalho()}.
        \end{itemize}
    \end{exampleblock}
\end{frame}

% Questao 17
\begin{frame}{Questão 17: Função Sem Parâmetro e Sem Retorno}
    \begin{block}{Quesito: Menu de Ajuda ao Operador}
        Implemente a função \texttt{exibir\_menu\_ajuda()} que mostre opções de teclas de atalho de um terminal (\texttt{[S] Salvar}, \texttt{[C] Carregar}, \texttt{[Q] Sair}). \\
        O programa principal deve chamar a função quando o usuário digitar \texttt{"ajuda"}.
    \end{block}

    \begin{exampleblock}{Ajuda de Resolução (Como Pensar)}
        \begin{itemize}
            \item Escreva a rotina contendo apenas saídas de tela organizadas.
            \item No fluxo principal, compare a entrada lida: \texttt{if opcao == "ajuda":}.
            \item Dentro do \texttt{if}, invoque a função pelo seu nome.
        \end{itemize}
    \end{exampleblock}
\end{frame}

% Questao 18
\begin{frame}{Questão 18: Função Com Parâmetro e Sem Retorno}
    \begin{block}{Quesito: Notificador de Serviços de Infraestrutura}
        Crie a função \texttt{notificar\_servico(nome, porta, status)} onde \texttt{status} é booleano. A função deve exibir:
        \begin{itemize}
            \item Se \texttt{True}: \texttt{"[OK] Serviço: <nome> | Porta: <porta> -> ATIVO"}.
            \item Se \texttt{False}: \texttt{"[FALHA] Serviço: <nome> | Porta: <porta> -> PARADO"}.
        \end{itemize}
    \end{block}

    \begin{exampleblock}{Ajuda de Resolução (Como Pensar)}
        \begin{itemize}
            \item Receba 3 argumentos: \texttt{def notificar\_servico(nome, porta, status):}.
            \item Avalie o parâmetro booleano com \texttt{if status: ... else: ...}.
            \item Exiba a mensagem formatada com f-string sem usar \texttt{return}.
        \end{itemize}
    \end{exampleblock}
\end{frame}

% Questao 19
\begin{frame}{Questão 19: Função Com Parâmetro e Sem Retorno}
    \begin{block}{Quesito: Divisor Visual Parametrizado}
        Crie a função \texttt{desenhar\_separador(simbolo, comprimento)} que receba um caractere textual e a quantidade de repetições horizontais. A função deve imprimir a linha gerada no console.
    \end{block}

    \begin{exampleblock}{Ajuda de Resolução (Como Pensar)}
        \begin{itemize}
            \item Em Python, multiplicar string por número repete o texto: \texttt{simbolo * comprimento}.
            \item Imprima a string resultante diretamente na tela.
            \item Teste com diferentes argumentos, como: \texttt{desenhar\_separador("-", 30)}.
        \end{itemize}
    \end{exampleblock}
\end{frame}

% Questao 20
\begin{frame}{Questão 20: Função Sem Parâmetro e Com Retorno}
    \begin{block}{Quesito: Leitor com Validação de Entrada Estrita}
        Implemente a função \texttt{ler\_inteiro\_positivo()} sem parâmetros.
        \begin{itemize}
            \item Solicite continuamente um número inteiro até que seja maior que zero;
            \item Assim que receber um valor estritamente positivo, devolva-o com \texttt{return}.
        \end{itemize}
    \end{block}

    \begin{exampleblock}{Ajuda de Resolução (Como Pensar)}
        \begin{itemize}
            \item Use um laço \texttt{while True:} dentro da função.
            \item Leia o número com \texttt{int(input())}.
            \item Se \texttt{valor > 0}, execute \texttt{return valor}.
            \item O \texttt{return} encerra o laço e a função no mesmo instante.
        \end{itemize}
    \end{exampleblock}
\end{frame}

% Questao 21
\begin{frame}{Questão 21: Função Sem Parâmetro e Com Retorno}
    \begin{block}{Quesito: Gerador de Código de Autenticação (OTP)}
        Escreva a função \texttt{gerar\_token\_acesso()} sem parâmetros que sorteie e retorne um código de $6$ dígitos inteiros (entre $100000$ e $999999$), simulando um token de segurança de dois fatores.
    \end{block}

    \begin{exampleblock}{Ajuda de Resolução (Como Pensar)}
        \begin{itemize}
            \item Utilize o módulo \texttt{random} do Python.
            \item Use a função \texttt{random.randint(100000, 999999)}.
            \item Retorne o valor sorteado com a instrução \texttt{return}.
            \item No programa principal, guarde o retorno em uma variável.
        \end{itemize}
    \end{exampleblock}
\end{frame}

% Questao 22
\begin{frame}{Questão 22: Função Com Parâmetro e Com Retorno}
    \begin{block}{Quesito: Conversor Multi-Escalas de Temperatura}
        Crie a função \texttt{converter\_celsius(temp\_c, escala)} onde escala é \texttt{"F"} ou \texttt{"K"}:
        \begin{itemize}
            \item Fórmulas: $F = C \times 1.8 + 32$ e $K = C + 273.15$.
            \item Calcule e retorne o resultado numérico; se a escala for inválida, retorne \texttt{None}.
        \end{itemize}
    \end{block}

    \begin{exampleblock}{Ajuda de Resolução (Como Pensar)}
        \begin{itemize}
            \item Use \texttt{escala.upper()} para aceitar maiúsculas e minúsculas.
            \item Estruture um \texttt{if-elif-else} para checar \texttt{"F"} e \texttt{"K"}.
            \item Calcule a conversão e faça \texttt{return} do valor obtido.
        \end{itemize}
    \end{exampleblock}
\end{frame}

% Questao 23
\begin{frame}{Questão 23: Função Com Parâmetro e Com Retorno}
    \begin{block}{Quesito: Predicado Booleano de Número Primo}
        Crie a função \texttt{eh\_primo(numero)} que receba um inteiro positivo e retorne \texttt{True} se for primo, ou \texttt{False} caso não seja.
    \end{block}

    \begin{exampleblock}{Ajuda de Resolução (Como Pensar)}
        \begin{itemize}
            \item Se o número for $\le 1$, retorne imediatamente \texttt{False}.
            \item Faça um laço \texttt{for d in range(2, int(numero**0.5) + 1):}.
            \item Se encontrar divisor com resto zero (\texttt{numero \% d == 0}), retorne \texttt{False}.
            \item Se o laço terminar sem divisores, retorne \texttt{True}.
        \end{itemize}
    \end{exampleblock}
\end{frame}

% Questao 24
\begin{frame}{Questão 24: Função Com Parâmetro e Retorno Múltiplo}
    \begin{block}{Quesito: Localizador de Extremos e Posição em Vetor}
        Crie a função \texttt{analisar\_pico(leituras)} que receba uma lista de números reais. \\
        A função deve retornar dois valores simultaneamente: o maior valor encontrado e o índice da primeira ocorrência desse maior valor.
    \end{block}

    \begin{exampleblock}{Ajuda de Resolução (Como Pensar)}
        \begin{itemize}
            \item Inicialize \texttt{maior = leituras[0]} e \texttt{posicao = 0}.
            \item Percorra a lista com índices usando \texttt{for i in range(1, len(leituras)):}.
            \item Se \texttt{leituras[i] > maior}, atualize ambas as variáveis.
            \item Devolva as duas variáveis com \texttt{return maior, posicao}.
        \end{itemize}
    \end{exampleblock}
\end{frame}

% Questao 25
\begin{frame}{Questão 25: Função Com Parâmetro e Com Retorno de Vetor}
    \begin{block}{Quesito: Filtro e Higienização de Conjunto de Dados}
        Crie a função \texttt{filtrar\_acima\_limiar(valores, limiar)} que receba uma lista e um número de corte. Retorne uma \textbf{nova lista} apenas com os elementos estritamente maiores que o limiar.
    \end{block}

    \begin{exampleblock}{Ajuda de Resolução (Como Pensar)}
        \begin{itemize}
            \item Inicialize uma lista vazia: \texttt{resultado = []}.
            \item Itere sobre os itens com \texttt{for item in valores:}.
            \item Se \texttt{item > limiar:}, adicione à nova lista com \texttt{resultado.append(item)}.
            \item Finalize retornando a nova lista: \texttt{return resultado}.
        \end{itemize}
    \end{exampleblock}
\end{frame}

\section{Resumo e Orientações Finais}

% Slide de Matriz de Funcoes
\begin{frame}{Matriz de Decisão: As Quatro Formas de Funções}
    \begin{table}
        \centering
        \scriptsize
        \begin{tabular}{lccl}
            \toprule
            \textbf{Classificação} & \textbf{Entrada} & \textbf{Saída} & \textbf{Exemplo Típico} \\
            \midrule
            Sem Parâm. e Sem Retorno & Não & Não & Banners e menus no console \\
            Com Parâm. e Sem Retorno & Sim & Não & Mensagens de log e alertas \\
            Sem Parâm. e Com Retorno & Não & Sim & Leituras com validação, tokens \\
            Com Parâm. e Com Retorno & Sim & Sim & Cálculos, filtros e buscas \\
            \bottomrule
        \end{tabular}
        \caption{Matriz dos Quatro Tipos de Funções em Programação.}
    \end{table}

    \vspace{0.2cm}
    \begin{block}{Regra Fundamental}
        \begin{itemize}
            \item \textbf{Entrada:} Recebida pelos parâmetros entre os parênteses \texttt{(...)}.
            \item \textbf{Saída:} Devolvida ao código chamador por meio da instrução \texttt{return}.
        \end{itemize}
    \end{block}
\end{frame}

% Slide de Roteiro de Resolucao de Problemas
\begin{frame}{Checklist para Resolução de Algoritmos}
    \begin{columns}
        \begin{column}{0.5\textwidth}
            \begin{block}{1. Análise Prévia}
                \begin{itemize}
                    \item Identifique as entradas necessárias.
                    \item Defina as saídas esperadas.
                    \item Escolha os tipos (\texttt{int}, \texttt{float}, \texttt{list}).
                \end{itemize}
            \end{block}
        \end{column}
        \begin{column}{0.5\textwidth}
            \begin{block}{2. Construção e Teste}
                \begin{itemize}
                    \item Isole a lógica em blocos ou funções.
                    \item Realize o teste de mesa no papel.
                    \item Teste valores de borda ($0$, negativos).
                \end{itemize}
            \end{block}
        \end{column}
    \end{columns}

    \vspace{0.3cm}
    \begin{center}
        \textit{"Primeiro resolva o problema. Depois escreva o código."}
    \end{center}
\end{frame}

% Slide de Fechamento
\begin{frame}{Encerramento e Bons Estudos!}
    \begin{center}
        \LARGE \textbf{Excelente trabalho na revisão!}\\
        \vspace{0.3cm}
        \normalsize Pratique a implementação de cada um dos 25 quesitos.\\
        \vspace{0.5cm}
        \textbf{Prof. Reginaldo Fernandes}\\
        \footnotesize \texttt{reginaldo.fernandes@ifce.edu.br}\\
        \vspace{0.3cm}
        \textbf{IFCE Campus Tauá}\\
        \small Curso Técnico em Redes de Computadores
    \end{center}
\end{frame}


\end{document}
